\section{Giai đoạn 1: Truy xuất}

\subsection{Thiết kế thực nghiệm}

\begin{itemize}
  \item \textbf{Mục tiêu:} chọn cấu hình đưa chunk chứa bằng chứng lên đầu danh
  sách trước khi đánh giá generation.
  \item \textbf{Phạm vi:} chỉ đánh giá retrieval offline; không gọi generator,
  không chấm RAGAS và không phát sinh chi phí API sinh văn bản.
  \item \textbf{Gold label:} đối chiếu kết quả với \path{relevant_chunk_ids}
  được ánh xạ từ evidence span. Một chunk chỉ đúng khi ID của nó thuộc tập gold.
\end{itemize}

\noindent\textbf{Thiết lập cố định}
\begin{itemize}
  \item Development: 281 câu resolved thuộc 50 bài; corpus gồm 11.064 bài.
  \item Chunking: recursive, kích thước 512 token và overlap 64 token.
  \item Resolved là biến thể chính vì câu hỏi có thể hiểu độc lập. Original chỉ
  dùng để phân tích ảnh hưởng của câu hỏi thiếu chủ thể.
\end{itemize}

\noindent\textbf{Quy trình chọn cấu hình}

\begin{enumerate}
  \item Sàng lọc riêng dense và sparse retriever, không rerank.
  \item So sánh dense tốt nhất, sparse tốt nhất và hybrid RRF với ba chế độ
        reranker.
  \item Khóa pipeline rồi chạy đúng một lần trên 871 câu thuộc 150 bài chưa
        thấy.
\end{enumerate}

\noindent\textbf{Tiêu chí đánh giá}
\begin{itemize}
  \item Ưu tiên MRR@5 vì generator có lợi nhất khi bằng chứng đầu tiên xuất hiện
  sớm.
  \item Khi MRR gần nhau, lần lượt xét nDCG@5, Hit@5 và latency.
  \item Dùng Recall@5 để kiểm tra độ phủ khi một câu có nhiều gold chunks.
  \item Mọi run phải đạt coverage 100\%; question IDs, seed và cách chấm được
  giữ cố định trong cùng một vòng.
\end{itemize}

\subsection{Sàng lọc retriever}

\begin{itemize}
  \item \textbf{Dense:} all-MiniLM-L6-v2, bge-small-en-v1.5,
  bge-large-en-v1.5 và e5-base-v2.
  \item \textbf{Sparse:} BM25 Okapi, BM25+, BM25 có
  stemming/stopword normalization và learned sparse BGE-M3.
  \item \textbf{Điều kiện chung:} mỗi phương pháp trả tối đa 10 candidates;
  chưa dùng cross-encoder.
\end{itemize}

\begin{table}[H]
  \centering
  \caption{Kết quả sàng lọc retriever trên 281 câu resolved.}
  \label{tab:phase1-retriever-screening}
  \small
  \begin{tabular}{lrrrr}
    \toprule
    \textbf{Retriever} & \textbf{MRR@5} & \textbf{nDCG@5} &
    \textbf{Hit@5} & \textbf{P50 (ms)} \\
    \midrule
    BGE-M3 learned sparse          & \textbf{0,8059} & \textbf{0,8317} & \textbf{0,9181} & 76,7 \\
    BM25 Okapi + stemming          & 0,7815 & 0,8123 & 0,9146 & 56,1 \\
    BM25+                          & 0,6881 & 0,7206 & 0,8221 & 101,1 \\
    BM25 Okapi                     & 0,6732 & 0,7087 & 0,8185 & 111,8 \\
    Dense: e5-base-v2              & \textbf{0,6316} & \textbf{0,6684} & \textbf{0,7865} & 15,5 \\
    Dense: bge-small-en-v1.5       & 0,6285 & 0,6589 & 0,7651 & 16,0 \\
    Dense: bge-large-en-v1.5       & 0,6252 & 0,6536 & 0,7544 & 27,4 \\
    Dense: all-MiniLM-L6-v2        & 0,4851 & 0,5165 & 0,6299 & 10,6 \\
    \bottomrule
  \end{tabular}
\end{table}

\noindent\textbf{Kết quả và phạm vi kết luận}
\begin{itemize}
  \item Trong các mô hình và dữ liệu được khảo sát, nhánh sparse đạt kết quả
  retrieval cao hơn nhánh dense.
  \item BGE-M3 được chọn làm sparse representative vì đứng đầu theo quy tắc
  MRR@5, sau đó là nDCG@5. Tuy nhiên, so với BM25-stemmed trên tập resolved,
  CI95 của hiệu nDCG@5 là $[-0,0125;\,0,0508]$; chưa đủ bằng chứng kết luận một
  sparse model tốt hơn model còn lại.
  \item e5-base-v2 được chọn làm dense representative. Chênh lệch giữa các
  dense models nhỏ và chưa đủ để thiết lập một thứ hạng tổng quát.
  \item BGE-M3 chậm hơn e5-base-v2; P95 vẫn dưới 0,1 giây trong thử nghiệm.
\end{itemize}

\subsection{Sàng lọc chiến lược retrieval và reranker}

\begin{itemize}
  \item \textbf{Retriever:} e5-base-v2 dense, BGE-M3 sparse và hybrid RRF giữa
  hai nhánh.
  \item \textbf{Reranker:} không rerank, cross-encoder MiniLM-L6 và
  BGE-reranker-large.
  \item \textbf{Thiết lập:} retriever lấy 20 candidates; reranker sắp lại và
  trả top 5 để tính metric.
\end{itemize}

\begin{table}[H]
  \centering
  \caption{MRR@5 và nDCG@5 theo retriever và reranker trên tập development.}
  \label{tab:phase1-reranker-matrix}
  \small
  \setlength{\tabcolsep}{4.5pt}
  \begin{tabular}{lrrr|rrr}
    \toprule
    & \multicolumn{3}{c|}{\textbf{MRR@5}}
    & \multicolumn{3}{c}{\textbf{nDCG@5}} \\
    \textbf{Retriever} & \textbf{Không} & \textbf{MiniLM} & \textbf{BGE-large}
    & \textbf{Không} & \textbf{MiniLM} & \textbf{BGE-large} \\
    \midrule
    BGE-M3 sparse & 0,8059 & 0,8387 & \textbf{0,8797} & 0,8317 & 0,8642 & \textbf{0,8976} \\
    e5-base dense & 0,6316 & 0,7783 & 0,7997 & 0,6684 & 0,7994 & 0,8172 \\
    Hybrid RRF    & 0,7192 & 0,7938 & 0,8203 & 0,7474 & 0,8164 & 0,8405 \\
    \bottomrule
  \end{tabular}
\end{table}

\begin{table}[H]
  \centering
  \caption{Độ trễ P50 của ma trận Phase 1 trên cùng môi trường chạy.}
  \label{tab:phase1-reranker-latency}
  \small
  \begin{tabular}{lrrr}
    \toprule
    \textbf{Retriever} & \textbf{Không rerank} & \textbf{MiniLM} &
    \textbf{BGE-large} \\
    \midrule
    BGE-M3 sparse & 66,1 ms & 163,9 ms & 510,1 ms \\
    e5-base dense & 17,7 ms & 127,3 ms & 540,5 ms \\
    Hybrid RRF    & 90,4 ms & 193,3 ms & 609,8 ms \\
    \bottomrule
  \end{tabular}
\end{table}

\noindent\textbf{Kết quả}
\begin{itemize}
  \item Trên BGE-M3, MiniLM tăng MRR@5 từ 0,8059 lên 0,8387; BGE-large đạt
  0,8797 nhưng P50 tăng từ 163,9 ms lên 510,1 ms.
  \item BGE-M3 + BGE-large đứng đầu ma trận: MRR@5 0,8797; nDCG@5 0,8976;
  Hit@5 0,9573; Recall@5 0,9555. Vì vậy, cấu hình này được chọn.
  \item Hybrid RRF không vượt BGE-M3 sparse ở cả ba chế độ reranker; fusion
  không đem lại lợi ích quan sát được trong thiết lập này.
\end{itemize}

\subsection{Final-test và cấu hình retrieval khóa}

\begin{itemize}
  \item \textbf{Cấu hình khóa trước final-test:} recursive chunking 512/64;
  BGE-M3 learned sparse lấy top 20; BGE-large cross-encoder trả top 5.
  \item \textbf{Dữ liệu:} 871 câu resolved thuộc 150 bài chưa dùng để chọn cấu
  hình.
  \item \textbf{Phạm vi:} chỉ chạy retrieval; không gọi generator. Cả 871/871
  câu đều có bản ghi thành công.
\end{itemize}

\begin{table}[H]
  \centering
  \caption{Kết quả của pipeline retrieval đã khóa.}
  \label{tab:phase1-final-test}
  \small
  \begin{tabular}{lrr}
    \toprule
    \textbf{Metric} & \textbf{Development (281 câu)} &
    \textbf{Final-test (871 câu)} \\
    \midrule
    MRR@5       & 0,8797 & 0,8112 \\
    nDCG@5      & 0,8976 & 0,8313 \\
    Hit@5       & 0,9573 & 0,8978 \\
    Recall@5    & 0,9555 & 0,8955 \\
    P50 latency & 512,7 ms & 550,0 ms \\
    P95 latency & 531,0 ms & 567,0 ms \\
    Coverage    & 100\% & 100\% \\
    \bottomrule
  \end{tabular}
\end{table}

\noindent\textbf{Kết quả}
\begin{itemize}
  \item 782/871 câu có ít nhất một gold chunk trong top 5; 89 câu không đạt
  Hit@5.
  \item So với BGE-M3 trước rerank, BGE-large tăng MRR@5 từ 0,7269 lên 0,8112,
  nDCG@5 từ 0,7555 lên 0,8313 và Hit@5 từ 0,8519 lên 0,8978.
  \item Mức tăng nDCG@5 là 0,0758, với paired CI95
  $[0,0504;\,0,1017]$. Khoảng này không chứa 0, nên hiệu quả rerank vẫn xuất
  hiện trên các bài chưa thấy.
\end{itemize}

\noindent\textbf{Diễn giải và quyết định}
\begin{itemize}
  \item Bốn metric chính đều thấp hơn development; điểm development đã lạc
  quan hơn so với dữ liệu chưa thấy.
  \item Coverage vẫn đạt 100\%; pipeline tìm được gold evidence trong top 5 cho
  gần 90\% câu hỏi.
  \item Khóa BGE-M3 top 20 + BGE-large top 5 làm retrieval đầu vào cho Phase 2.
  Final-test không được dùng để chọn lại model.
\end{itemize}
